\documentclass[10pt,a4paper]{amsart}
\usepackage[margin=19mm]{geometry}
\usepackage{amsmath,amssymb,mathtools}
\usepackage[scr=rsfs,cal=euler]{mathalfa}
\usepackage{xfrac}
\usepackage[unicode,hidelinks,bookmarks=false]{hyperref}
\newcommand{\ie}{i.e.\ }
\newcommand{\cf}{cf.\ }
\newcommand{\resp}{respectively\ }
\newcommand{\Subsection}{\subsection}
\providecommand{\nobreakdash}{\nobreak\hbox{-}}
\setlength{\emergencystretch}{2em}
\multlinegap0pt
\newtheorem{definition}{Definition}
\newtheorem{theorem}{Theorem}
\newtheorem{lemma}{Lemma}
\newtheorem{example}{Example}
\newtheorem{corollary}{Corollary}
\newtheorem{remark}{Remark}
\begin{document}
\title[Spectral Collapse Under Geometric Alignment of Extreme Event]{Spectral Collapse Under Geometric Alignment of Extreme Events}
\author{Jokubas Petkevicius}
\date{}
\maketitle
\noindent\emph{Source and licence.} Spectral Collapse Under Geometric Alignment of Extreme Events. Version arXiv:2606.25810v1 (CC BY 4.0). This material is distributed under the Creative Commons Attribution 4.0 International licence, http://creativecommons.org/licenses/by/4.0/, as recorded in the arXiv OAI licence metadata for this exact version and shown on the abstract page https://arxiv.org/abs/2606.25810v1. Full rights notice: licenses/arxiv-2606.25810-CC-BY-4.0.txt.\par\smallskip
\noindent\textit{Editorial scope.} Counted unit: the original \texttt{theorem} environment \texttt{-} at source lines 357--395 together with its complete proof at lines 397--454; the delivered document keeps source lines 353--455 so that every referenced label and every citation resolves inside it. Native editable source is the paper's own concatenated TeX; the AMS wrapper is editorial formatting only. No publisher class, upstream script or unrelated artwork is redistributed.
\section{Main Theorem: Spectral Collapse of the Full System}

We now combine the geometric alignment result with background negligibility to obtain collapse of the full system.


\begin{theorem}[Spectral Collapse Under Aligned Jumps]
Let

\[
Q_n = B_n + J_n, \quad J_n = \sum_{k=1}^{m_n} \theta_{n,k} v_{n,k} v_{n,k}^\top,
\]

where:
\begin{itemize}
\item \(B_n\) is symmetric positive semidefinite
\item \(\|v_{n,k}\| = 1\), \(\theta_{n,k} \ge 0\)
\end{itemize}

Assume the following asymptotic conditions as \(n \to \infty\):

\begin{enumerate}
\item \textbf{Geometric alignment of jumps:} There exists a unit vector \(u_n\) such that
\[
\sum_{k=1}^{m_n} \theta_{n,k} \left(1 - \langle v_{n,k}, u_n \rangle^2 \right) = o\left(\sum_{k=1}^{m_n} \theta_{n,k}\right).
\]

\item \textbf{Trace dominance:} \(\operatorname{tr}(B_n) = o(\operatorname{tr}(J_n))\)

\item \textbf{Second moment dominance:} \(\operatorname{tr}(B_n^2) = o(\lambda_1(J_n)^2)\)
\end{enumerate}

Then:

\[
\frac{\lambda_1(Q_n)}{\operatorname{tr}(Q_n)} \to 1,
\]

and

\[
\operatorname{erank}(Q_n) \to 1.
\]

\end{theorem}

\begin{proof}
By Theorem 1, condition (1) implies
\[
\frac{\lambda_1(J_n)}{\operatorname{tr}(J_n)} \to 1.
\]

\textbf{Step 1: Lower bound on \(\lambda_1(Q_n)\).}
Let \(u_n\) be the vector from condition (1). Then
\[
\lambda_1(Q_n) \ge u_n^\top Q_n u_n = u_n^\top B_n u_n + u_n^\top J_n u_n.
\]
We have \(u_n^\top J_n u_n = \operatorname{tr}(J_n) - o(\operatorname{tr}(J_n))\) and \(u_n^\top B_n u_n \ge -\|B_n\|_{\mathrm{op}}\). Hence
\[
\lambda_1(Q_n) \ge \operatorname{tr}(J_n) - o(\operatorname{tr}(J_n)) - \|B_n\|_{\mathrm{op}}.
\]

\textbf{Step 2: Relating \(\|B_n\|_{\mathrm{op}}\) to \(\operatorname{tr}(J_n)\).}
From condition (3), \(\|B_n\|_{\mathrm{op}}^2 \le \operatorname{tr}(B_n^2) = o(\lambda_1(J_n)^2)\). Since \(\lambda_1(J_n) \sim \operatorname{tr}(J_n)\), we have \(\|B_n\|_{\mathrm{op}} = o(\operatorname{tr}(J_n))\).

\textbf{Step 3: Upper bound on \(\lambda_k(Q_n)\) for \(k \ge 2\).}
By Weyl's inequality,
\[
\lambda_k(Q_n) \le \lambda_k(J_n) + \|B_n\|_{\mathrm{op}}.
\]
For \(k \ge 2\), \(\lambda_k(J_n) \le \operatorname{tr}(J_n) - \lambda_1(J_n) = o(\operatorname{tr}(J_n))\). Hence
\[
\lambda_k(Q_n) = o(\operatorname{tr}(J_n)) + o(\operatorname{tr}(J_n)) = o(\operatorname{tr}(J_n)).
\]

\textbf{Step 4: Ratio convergence.}
We have
\[
\operatorname{tr}(Q_n) = \operatorname{tr}(B_n) + \operatorname{tr}(J_n) = \operatorname{tr}(J_n) + o(\operatorname{tr}(J_n)).
\]
From Step 1,
\[
\lambda_1(Q_n) \ge \operatorname{tr}(J_n) - o(\operatorname{tr}(J_n)).
\]
From Step 3,
\[
\operatorname{tr}(Q_n) = \lambda_1(Q_n) + \sum_{k=2}^n \lambda_k(Q_n) = \lambda_1(Q_n) + o(\operatorname{tr}(J_n)).
\]
Thus \(\lambda_1(Q_n) = \operatorname{tr}(J_n) + o(\operatorname{tr}(J_n)) = \operatorname{tr}(Q_n) + o(\operatorname{tr}(Q_n))\), so
\[
\frac{\lambda_1(Q_n)}{\operatorname{tr}(Q_n)} \to 1.
\]

\textbf{Step 5: Effective rank collapse.}
We compute
\[
\operatorname{tr}(Q_n^2) = \lambda_1(Q_n)^2 + \sum_{k=2}^n \lambda_k(Q_n)^2 = \lambda_1(Q_n)^2 + o(\operatorname{tr}(J_n)^2) = \lambda_1(Q_n)^2 + o(\lambda_1(Q_n)^2).
\]
Hence
\[
\operatorname{erank}(Q_n) = \frac{(\operatorname{tr} Q_n)^2}{\operatorname{tr}(Q_n^2)} = \frac{\lambda_1(Q_n)^2 + o(\lambda_1(Q_n)^2)}{\lambda_1(Q_n)^2 + o(\lambda_1(Q_n)^2)} \to 1.
\]

\end{proof}


\end{document}
